Patch-based tokenisation has become the default input stage of time-series foundation models, adopted because it shortens sequences and makes attention affordable. The price is assumed to be resolution: fine detail traded for length, with the sampling theorem deciding what survives. This report argues the price is of a different kind. Where the period of a signal is commensurate with the patch geometry, consecutive patches repeat, and a frequency the sampler represented perfectly becomes poorly observable to the model. Nyquist has nothing to say about that loss, because the tokeniser causes it rather than undersampling.

The distinction matters wherever such a model sits in a monitoring or control loop. A blind spot that follows from the configuration rather than from the data is invisible to an operator who never derived it, stable enough to be relied on by anyone who did, and unreachable by a detector placed downstream of the tokeniser. It is an engineering defect and a security property at once, and is treated here as both.

\autoref{sec:introduction} gives the background. \autoref{sec:one} and \autoref{sec:two} formalise the tokenisation operator and derive the locking condition. \autoref{sec:three} and \autoref{sec:four} fix the semantic layer and state the hypotheses the design must be able to refute. \autoref{sec:five} and \autoref{sec:six} describe the signals and the experiments, \autoref{sec:seven} what they return, \autoref{sec:eight} one mitigation, \autoref{sec:nine} the agentic deployment and \autoref{sec:ten} the risk assessment. Six appendices hold the ontology, the Bayesian specifications, the generators, the retraining recipe and the exploratory sweep.
